
\documentclass[12pt,a4paper]{amsart}

\usepackage{amsmath,amssymb,amsthm}
\usepackage{mathtools}
\usepackage{enumitem}
\usepackage{hyperref}

\newtheorem{theorem}{Theorem}[section]
\newtheorem{lemma}[theorem]{Lemma}
\newtheorem{proposition}[theorem]{Proposition}
\newtheorem{corollary}[theorem]{Corollary}
\theoremstyle{definition}
\newtheorem{definition}[theorem]{Definition}
\newtheorem{assumption}[theorem]{Assumption}
\newtheorem{problem}[theorem]{Open Problem}
\theoremstyle{remark}
\newtheorem{remark}[theorem]{Remark}

\newcommand{\Rig}{\textup{[Rig]}}
\newcommand{\Cond}{\textup{[Cond]}}
\newcommand{\Open}{\textup{[Open]}}
\newcommand{\Skel}{\textup{[Skeleton]}}

\begin{document}

\title[A$_2$-Stability for PSWF Amplitudes]{%
  A$_2$-Stability for Prolate Spheroidal Wave Function Amplitudes\\
  in the Semiclassical Limit:\\
  Proof Architecture Skeleton}

\author{[Author]}

\begin{abstract}
This is the \emph{Level~0--1 proof architecture skeleton} for Paper~XVII.
It fixes the logical structure and identifies the critical bottlenecks
before full LaTeX elaboration.

\medskip
\noindent\textbf{Objective.}
Verify Hypothesis~\textup{hyp:rigidity} from Paper~XVI
for PSWF amplitudes:
the PSWF amplitude satisfies Assumption~A2 of Paper~XVI
(conditions (A1)--(A3)) uniformly as $h := c^{-1} \to 0$.

\medskip
\noindent\textbf{Central new object.}
The \emph{Olver--FIO Lift Lemma}
(Section~\ref{sec:lift}):
uniform Olver asymptotics define a Lagrangian distribution
whose associated canonical relation carries a single A$_2$-fold,
and whose CFU Jacobian admits a uniform lower bound in $h$.
This is not available in the standard literature
and constitutes the real analytic contribution of Paper~XVII.

\medskip
\noindent\textbf{Strategy.}
Primary: Olver turning-point analysis (Strategy~B).
Secondary: FIO reinterpretation of Olver expansions (Strategy~A).
Auxiliary: Sturm--Liouville spectral theory for global consistency.
\end{abstract}

\maketitle

\tableofcontents

\section{Input from Paper~XVI}
\label{sec:input}

The following are used as \emph{black boxes} (not reproved here):

\begin{itemize}
\item \textbf{thm:CFU-geom} (XVI):
  A$_2$-fold phase $\Rightarrow$ Airy normal form
  via local canonical transformation $\kappa$.
\item \textbf{thm:CFU-symbol} (XVI):
  $I^{-1/4,0}(\Lambda_0,\Lambda_1) \to S^{-1/2}_{1,0}(h)$
  under $\kappa$, uniform in $h$.
\item \textbf{Assumption~A2} (XVI):
  Three conditions required for CFU reduction:
  \begin{enumerate}[{\rm(A1)}]
  \item Single nondegenerate fold:
    $c_3(h) \ge \delta_0 > 0$ uniformly.
  \item CFU admissibility:
    $|\det(dt/d\beta)| \ge \delta_1 > 0$ uniformly in $h$
    (the hard bottleneck).
  \item Unique fold on $\mathrm{supp}\,\chi$.
  \end{enumerate}
\item \textbf{hyp:rigidity} (XVI, conjectural):
  Airy class $\Rightarrow$ A$_2$-fold uniform in $h$
  (the inverse direction to be verified here for PSWF).
\end{itemize}

\section{PSWF Setup}
\label{sec:pswf}

\begin{definition}[PSWF as semiclassical eigenfamily]\label{def:pswf}
The prolate spheroidal wave functions $\psi_n(x;c)$,
$x \in [-1,1]$, are eigenfunctions of the
differential operator
\[
  \mathcal{L}_c := -\frac{d}{dx}
  \bigl[(1-x^2)\frac{d}{dx}\bigr]
  + c^2 x^2,
\]
with eigenvalues $\chi_n(c)$.
Setting $h := c^{-1}$, the operator becomes
\[
  h^2 \mathcal{L}_{h^{-1}}
  = -h^2\frac{d}{dx}\bigl[(1-x^2)\frac{d}{dx}\bigr]
  + x^2,
\]
a semiclassical Schr\"odinger operator with
potential $V(x) = x^2/(1-x^2)$ and
semiclassical parameter $h \to 0$.
The turning points (where $V(x) = \chi_n h^2$)
are the loci of A$_2$-fold singularities.
\end{definition}

\begin{remark}[Turning point structure]
\label{rem:turning}
For fixed mode $n$ and $h \to 0$,
the eigenvalue $\chi_n(c) \sim c^2 \lambda_n$
for some $\lambda_n \in (0,1)$.
The turning points $x_\pm(h)$ satisfy
$V(x_\pm) = \chi_n h^2$
and converge to the classical turning points
as $h \to 0$.
Each turning point is expected to be an
A$_2$-fold (Olver \cite{Olver1974}, Ch.~10--11);
verification of uniformity is the task of
Section~\ref{sec:lift}.
\end{remark}

\section{Core Strategy: Olver Turning-Point Analysis}
\label{sec:olver}

\subsection{Olver's uniform asymptotic scheme}

Olver \cite{Olver1974} (Ch.~10--11) provides,
for second-order ODEs with a simple turning point:
\begin{enumerate}[(i)]
\item A \emph{uniform asymptotic expansion}
  matching bulk (WKB) and turning-point (Airy) regimes.
\item \emph{Controlled error bounds}
  uniform in $h$ on compact subsets avoiding endpoints.
\item A local change of variables $x \mapsto t(x;h)$
  reducing the ODE phase to
  $\frac{2}{3}(-t)^{3/2}$ (classical Airy scaling).
\end{enumerate}

\begin{remark}[What Olver does \emph{not} give]
\label{rem:olver-gap}
Olver guarantees the \emph{existence} of the
change of variables $x \mapsto t(x;h)$
and the asymptotic structure of eigenfunctions,
but does \emph{not} directly provide:
\begin{itemize}
\item a lower bound $|\det(dt/d\beta)| \ge \delta_1 > 0$
  uniform in $h$ (this is Assumption (A2) of XVI);
\item a microlocal canonical relation framework;
\item the FIO-class membership of the amplitude.
\end{itemize}
Bridging this gap is the content of
Lemma~\ref{lem:olver-fio-lift} below.
\end{remark}

\subsection{FIO reinterpretation of Olver expansions}

The Olver expansion can be read as defining
a semiclassical Lagrangian distribution.
Specifically, the PSWF admits (schematically)
\[
  \psi_n(x;h)
  \sim h^{-1/4}\,
  a(x;h)\,
  \mathrm{Ai}\bigl(h^{-2/3}\,\phi(x;h)\bigr)
  + \text{(oscillatory tail)},
\]
where $\phi(x;h) \to \phi_0(x)$ as $h \to 0$
and $\phi_0(x_0) = 0$ at the turning point $x_0$.
The amplitude $a(x;h)$ carries the
$I^{m,l}$-structure of Paper~XVI.
Making this identification rigorous is the
central task.

\section{The Olver--FIO Lift Lemma}
\label{sec:lift}

This section contains the central new object of Paper~XVII.

\begin{lemma}[Olver--FIO Lift, \Skel]
\label{lem:olver-fio-lift}
\emph{(Skeleton formulation; proof to be elaborated.)}
Let $\psi_n(\,\cdot\,;h)$ be the PSWF family
(Definition~\ref{def:pswf}).
Under suitable uniformity conditions on the
eigenvalue branch $\chi_n(c)$
(to be made precise):
\begin{enumerate}[{\rm(i)}]
\item \textit{(Lagrangian distribution.)}
  $\psi_n(\,\cdot\,;h)$ is microlocally equivalent,
  modulo $S^{-\infty}$ uniformly in $h$,
  to a Lagrangian distribution associated to a
  canonical relation $\mathcal{C}_n \subset
  T^*[-1,1] \times T^*\mathbb{R}$.
\item \textit{(A$_2$-fold.)}
  $\mathcal{C}_n$ has a single A$_2$-fold
  at the turning point $(x_0,\xi_0)$
  (Assumption~(A1) and~(A3) of XVI),
  with $c_3(h) \ge \delta_0 > 0$ uniformly in $h$.
\item \textit{(Uniform Jacobian bound.)}
  The CFU change of variables
  $\beta \mapsto t(\beta,u;h)$ extracted from the
  Olver reduction satisfies
  $|\det(dt/d\beta)| \ge \delta_1 > 0$
  uniformly in $h$
  (Assumption~(A2) of XVI).
\end{enumerate}
Consequently, the PSWF amplitude satisfies
Assumption~A2 of Paper~XVI uniformly in $h \to 0$,
and \textup{hyp:rigidity} holds for the PSWF family.
\end{lemma}

\begin{remark}[Proof strategy for Lemma~\ref{lem:olver-fio-lift}]
\label{rem:lift-strategy}
\begin{enumerate}[{\rm(a)}]
\item Part~(i): identify Olver's asymptotic series
  as a Lagrangian distribution via the FIO calculus
  of H\"ormander \cite{Hormander1983} / Zworski \cite{Zworski2012}.
  The phase function is $\phi(x;h)$ from the
  Olver reduction.
\item Part~(ii): use Sturm--Liouville theory
  (unique simple turning point per mode)
  + Olver's local analysis to confirm A$_2$-type.
\item Part~(iii) is the hard step:
  requires controlling the Jacobian
  $|\partial t/\partial\beta|$ in the Olver
  change of variables uniformly in $h$.
  Candidate approach:
  explicit formula for $dt/d\beta$
  from the Olver phase and a lower bound
  via non-vanishing of $\phi'(x_0;h)$ uniformly.
  This requires a new estimate not in Olver \cite{Olver1974}.
\end{enumerate}
\end{remark}

\section{Critical Bottlenecks}
\label{sec:bottlenecks}

\subsection{(A2) Uniform Jacobian bound: the hard core}

The only genuinely new analytic input required is:
\[
  \inf_{h \in (0,h_0],\; x \in \mathrm{supp}\,\chi}
  |\phi'(x;h)| \ge \delta > 0.
\]
This would imply the Jacobian bound in
Lemma~\ref{lem:olver-fio-lift}(iii)
via the implicit function theorem.
Candidate proof:
\begin{itemize}
\item Show $\phi'(x_0;h) \ne 0$ for all $h \in (0,h_0]$
  by differentiating the Olver phase equation.
\item Uniform control via compactness argument
  on the eigenvalue branch.
\end{itemize}
Status: \Open.

\subsection{(A1) Geometric stability of fold}

Follows from Olver's local analysis if
$c_3(h) = \frac{1}{6}\Theta'''(0;h) \ge \delta_0 > 0$.
Expected from the turning-point expansion;
uniform remainder control needed.
Status: \Cond\ (expected provable).

\subsection{(A3) Uniqueness of fold}

Follows from Sturm--Liouville theory:
for fixed mode $n$, exactly one simple turning point
per interval $(-1,1)$.
Status: \Rig\ (standard).

\section{Conjectural vs.\ Provable Split}
\label{sec:status}

\begin{center}
\begin{tabular}{lll}
\hline
Statement & Status & Source \\
\hline
(A3) unique fold & \Rig & Sturm--Liouville \\
Airy scaling locally & \Rig & Olver Ch.~11 \\
(A1) $c_3(h) \ge \delta_0$ & \Cond & Olver + uniformity \\
FIO-class membership & \Cond & new argument needed \\
(A2) Jacobian bound & \Open & core of XVII \\
\textup{hyp:rigidity} for PSWF & \Cond & follows from (A1--A3) \\
\hline
\end{tabular}
\end{center}

\section{Target Theorem of Paper~XVII}
\label{sec:target}

\begin{theorem}[PSWF A$_2$-stability, \Cond]
\label{thm:xvii-main}
\emph{(Target; proof contingent on
Lemma~\ref{lem:olver-fio-lift}(iii).)}
The PSWF amplitude $\psi_n(\,\cdot\,;h)$
satisfies Assumption~A2 of Paper~XVI
uniformly as $h \to 0$.
Consequently, the CFU reduction
(Theorems~thm:CFU-geom, thm:CFU-symbol of XVI)
applies to the PSWF family,
and \textup{hyp:rigidity} (XVI)
holds for PSWF.
\end{theorem}

\begin{remark}[What Paper~XVII does not claim]
Paper~XVII does not attempt to:
prove the full inverse microlocal conjecture
in general (this remains open);
give global (non-microlocal) bounds on PSWF;
resolve the PSWF completeness or
concentration-of-energy problem.
The scope is strictly:
verify Assumption~A2 of XVI for PSWF
via Olver + FIO reinterpretation.
\end{remark}

\begin{thebibliography}{9}
\bibitem{PaperXVI}
[Author], \textit{Paper~XVI}, 2026.
\bibitem{Olver1974}
F.~W.~J.~Olver,
\textit{Asymptotics and Special Functions},
Academic Press, 1974.
\bibitem{Zworski2012}
M.~Zworski,
\textit{Semiclassical Analysis}, AMS, 2012.
\bibitem{Hormander1983}
L.~H\"ormander,
\textit{The Analysis of Linear PDE~I},
Springer, 1983.
\bibitem{CFU1957}
C.~Chester, B.~Friedman, F.~Ursell,
Proc.\ Cambridge Philos.\ Soc.\
\textbf{53} (1957), 599--611.
\bibitem{MelroseU1979}
R.~Melrose, G.~Uhlmann,
Comm.\ Pure Appl.\ Math.\
\textbf{32} (1979), 483--519.
\bibitem{Slepian1961}
D.~Slepian, H.~O.~Pollak,
Bell Syst.\ Tech.\ J.\
\textbf{40} (1961), 43--63.
\end{thebibliography}

\end{document}
